\documentclass[11pt,a4paper]{article}
\usepackage[margin=2.5cm]{geometry}
% Shared by main.tex (the full report) and summary.tex (the one-pager).
\usepackage[T1]{fontenc}
\usepackage[utf8]{inputenc}
\usepackage{lmodern}
\usepackage{microtype}
\usepackage{amsmath}
\usepackage{booktabs}
\usepackage{xcolor}
\usepackage{graphicx}
\usepackage{caption}
\usepackage{subcaption}
\usepackage{enumitem}
\usepackage[hidelinks]{hyperref}

% Keep every figure and table inside the section that introduces it: floats
% may not drift past the next \section. Without this, a short section's floats
% pile up at the end of the document - inside the references.
\IfFileExists{placeins.sty}{\usepackage[section]{placeins}}{}
\providecommand{\FloatBarrier}{}

% Figures are drawn by pgfplots straight from data/*.csv. Overleaf has
% pgfplots; a minimal local TeX install may not, in which case each figure
% compiles to a labelled placeholder box instead of failing the build.
\newif\ifplots
\IfFileExists{pgfplots.sty}{%
  \usepackage{pgfplots}%
  \pgfplotsset{compat=1.18}%
  \usepgfplotslibrary{statistics}%
  \plotstrue
}{}
\newcommand{\plotfigure}[1]{%
  \ifplots\input{#1}%
  \else\fbox{\parbox[c][4cm][c]{0.9\linewidth}{\centering\ttfamily
        pgfplots not installed here -- this figure (#1) renders on Overleaf}}%
  \fi}

% Every number that comes from an experiment, generated by make_data.py.
% generated by make_data.py - do not edit by hand
\newcommand{\numAnalyzeShare}{4.4}
\newcommand{\numCoefConflict}{4232}
\newcommand{\numCoefScan}{2.1}
\newcommand{\numCoefSearch}{86.8}
\newcommand{\numCoefSkip}{6.0}
\newcommand{\numGapFlipFull}{1.51}
\newcommand{\numGapFlipNoRestart}{0.77}
\newcommand{\numGapFlipNoRestartNoMin}{1.08}
\newcommand{\numGapNoflipFull}{0.78}
\newcommand{\numGapNoflipNoRestart}{0.91}
\newcommand{\numGapNoflipNoRestartNoMin}{1.06}
\newcommand{\numModelMedianErr}{2.5}
\newcommand{\numModelRuns}{39}
\newcommand{\numPadFootprint}{20}
\newcommand{\numPadSlowdown}{23}
\newcommand{\numProofOverheadMax}{23}
\newcommand{\numProofOverheadMin}{11}
\newcommand{\numPropagateShare}{93.2}
\newcommand{\numSearchSharePhp}{49}
\newcommand{\numSearchShareRand}{33}


% Visible reminders in the compiled PDF. Remove the macro body (or the
% \TODO calls) once a section is written.
\newcommand{\TODO}[1]{\textcolor{red!70!black}{\textsf{[TODO: #1]}}}

% Instance names in running text.
\newcommand{\inst}[1]{\texttt{#1}}


\title{Where a CDCL SAT Solver Spends Its Time:\\
       Profiling, a Cycle-Level Cost Model, and a Controlled Comparison with MiniSat}
\author{Alperen Buğra Başar \\ \small Bilkent University \\ \small \texttt{bugra.basar@ug.bilkent.edu.tr}}
\date{{October 2026}}

\begin{document}
\maketitle

\begin{abstract}
I study the performance of a CDCL SAT solver written from scratch in C++17. A
sampling profile places half of its run time in one search loop, and four
per-operation costs, fitted over \numModelRuns{} runs, predict its total cycles
with a median error of \numModelMedianErr\%. A memory experiment that holds the
search fixed while growing the footprint \numPadFootprint{} times locates the
cost of leaving the L2 cache. A single run suggested the solver was 12 times
faster than MiniSat; shuffled copies and removing MiniSat's features one at a
time show that the real difference in conflicts is exactly its restarts and
clause minimization. Every answer is checked independently, including a DRAT
proof for every unsatisfiable result.
\end{abstract}

\section{The solver, and why its answers can be trusted}
\label{sec:solver}

This report studies the performance of a CDCL (conflict-driven clause
learning) SAT solver that I wrote from scratch in C++17, with no external
dependencies. It implements the core of a modern CDCL
solver~\cite{marquessilva1999grasp,moskewicz2001chaff,een2003minisat}: two
watched literals per clause, first-UIP clause learning, non-chronological
backjumping, VSIDS branching with phase saving, clause deletion by
LBD~\cite{audemard2009glucose}, and DRAT proof output. It does not yet have
restarts, clause minimization, blocking literals, a heap for branching or
arena clause storage; Section~\ref{sec:roadmap} uses the measurements to
predict what each of these would bring.

A fast wrong answer is worthless, so no result in this report relies on the
solver's own word. A satisfying assignment is easy to check: the test harness
re-reads the formula from disk and checks the assignment clause by clause. An
unsatisfiable answer has no such short certificate, so the solver writes a
DRAT proof, which \texttt{drat-trim}~\cite{wetzler2014drattrim} checks
independently. The regression suite runs both checks on 1{,}010 instances, and
a fuzzer compares 15{,}000 random formulas against exhaustive brute force,
also with clause deletion forced to run early; there are no mismatches. An
invariant checker guards the internal data structures, and it caught all six
bugs I injected on purpose to test it.

Verification has a price (Table~\ref{tab:proofcost}). Writing the proof costs
\numProofOverheadMin--\numProofOverheadMax\% more cycles, and the overhead
follows how much proof is written per unit of solving work (Spearman
$\rho = 0.95$); proof logging never changes the search. Checking takes
0.7--1.5 times the solving time on the pigeonhole and random instances, but
9--27 times on the fullN family, which I have not investigated.
\texttt{drat-trim} can also produce LRAT, the input format of the formally
verified checker cake\_lpr~\cite{tan2021cakelpr}, so even the checker would
not have to be trusted.

\begin{table}[ht]
\centering\small
\begin{tabular}{lrrrrr}
\toprule
instance & conflicts & logging overhead & proof (MB) & check (s) & check / solve \\
\midrule
% generated by make_data.py - do not edit by hand
\texttt{php9\_8} & 10360 & +19.5\% & 1.2 & 0.12 & 1.5 \\
\texttt{php10\_9} & 60416 & +15.2\% & 9.9 & 1.13 & 1.3 \\
\texttt{php11\_10} & 245609 & +12.8\% & 54.0 & 7.12 & 1.3 \\
\texttt{full15} & 16384 & +22.6\% & 1.0 & 0.57 & 9.2 \\
\texttt{full16} & 32768 & +22.4\% & 2.3 & 2.89 & 19.7 \\
\texttt{full17} & 65536 & +20.4\% & 5.1 & 9.90 & 27.3 \\
\texttt{rand3\_200\_s1} & 37502 & +14.1\% & 4.1 & 0.33 & 0.9 \\
\texttt{rand3\_250\_s2} & 93978 & +12.6\% & 12.8 & 0.99 & 0.8 \\
\texttt{rand3\_250\_s3} & 108163 & +11.3\% & 14.7 & 1.18 & 0.7 \\
\bottomrule

\end{tabular}
\caption{Cost of an independently checked UNSAT answer. Overhead is extra
cycles with proof logging on; every proof was verified by \texttt{drat-trim}.}
\label{tab:proofcost}
\end{table}

\section{Method}
\label{sec:method}

All measurements come from one machine, an Apple M3 Pro whose performance
cores run at 4.0\,GHz. A small pointer-chasing program measured what a load
costs on it: about 4 cycles from the L1 cache (128\,KB), 15--36 cycles from L2
(16\,MB) and 360--450 cycles from main memory, but only 4--8 cycles at any size
when memory is read in order, because the processor then fetches the next data
in advance. These numbers are the scale for every cost in this report.

Timings use the optimized build of the solver (\texttt{-O2}, assertions off),
and profiles use the same build with debug symbols added. Instead of seconds
I report processor cycles and instructions, read from \texttt{/usr/bin/time},
since cycles do not depend on the clock speed, and I keep the fastest of three
runs. The solver also counts its own operations. These counts are identical on
every run of the same input, so whenever a change should not affect the search
(proof output, memory padding, the counters themselves) I checked that they
stayed identical. The counters cost 3--4\% of the run time; my first version,
which updated memory on every clause visit, cost 10\%.

The solver is deterministic, but the same formula written differently
(variables renamed, signs flipped, clauses reordered) can need twice as many
conflicts, because the heuristics take a different path. I therefore compare
configurations on shuffled copies of an instance, generated from fixed seeds,
and in pairs: both configurations run on the same copy, so each copy's own
difficulty cancels out in their ratio. A sign test and a Wilcoxon signed-rank
test on the log ratios tell whether an effect is real, and a bootstrap 95\%
confidence interval for the median ratio gives its size.

Every experiment is a script that writes a CSV file recording the git commit,
build flags, machine and command, and every figure, table and number in this
report is generated from those files (Appendix~\ref{app:repro}). Three limits
apply throughout. The processor's hardware event counters were not available
(the installed Xcode was too old for the operating system), so statements
about caches and branches are inferences from several independent
measurements, not direct counts. There is a single machine. And all instances
have at most 250 variables, so most of them fit in the L2 cache.

\section{Where the time goes}
\label{sec:profile}

I profiled with \texttt{sample}, the macOS profiler, which records the running
instruction once per millisecond; a run of \inst{php11\_10} gives about
3{,}400 samples, so a share is accurate to about $\pm 1$ percentage point. On
that instance \texttt{propagate()} takes \numPropagateShare\% of the run and
\texttt{analyze()} \numAnalyzeShare\%; everything else is below 1\%. Cheap is
not the same as unimportant, though: clause deletion takes almost no time
itself, yet it made \inst{php11\_10} 4.4 times faster by keeping the clause
database small.

Inside \texttt{propagate()}, about half of the whole run goes into searching a
clause for a new literal to watch after one of its watched literals has become
false (Figure~\ref{fig:regions}), and one instruction alone receives 37\% of
all samples. A sample lands on the instruction the processor is \emph{waiting
at}, not on the one that caused the wait, an effect called skid. The hot
instruction is a cheap load from L1, reachable only by a jump from the check
that ends the search loop. That check depends on the current assignment and is
hard to predict, and after a wrong guess the processor restarts exactly there.
This points to branch mispredictions, a hypothesis I cannot confirm without
hardware counters.

The random instance supports it. Its learnt clauses are shorter, so a search
checks about 2 literals instead of 10: the loop body shrinks from 8\% to
1.5\% of the run, but the loop exit stays expensive, so the cost is paid per
search, not per literal. Moving a watch, a store to an essentially random
address, costs only 2.4\%, because the processor waits for loads but not for
stores.

\begin{figure}[ht]
\centering
\plotfigure{figures/fig-regions}
\caption{Share of all profile samples in each part of \texttt{propagate()}.
Long learnt clauses (pigeonhole) make the search for a new watch dominant;
short ones (random 3-SAT) shift weight to the fixed cost of each visit.}
\label{fig:regions}
\end{figure}

\section{What a clause visit does}
\label{sec:visits}

Looking at one clause that watches a literal that has just become false is a
\emph{clause visit}, the basic unit of work: \inst{php11\_10} performs 594
million of them. A visit is a \emph{skip} if the clause is already satisfied,
a \emph{move} if another non-false literal takes over the watch, \emph{unit}
if the last literal is forced true, and a \emph{conflict} if that literal is
false as well. Counting all four gives two exact checks, since they must add up
to the number of visits and the conflict visits must equal the conflicts; both
hold on every instance.

\begin{table}[ht]
\centering\small
\begin{tabular}{lrrrrr}
\toprule
instance & skip (\%) & moved (\%) & unit (\%) & conflict (\%) & literals per search \\
\midrule
% generated by make_data.py - do not edit by hand
\texttt{php10\_9} & 70.6 & 28.6 & 0.81 & 0.072 & 9.5 \\
\texttt{php11\_10} & 73.6 & 25.8 & 0.49 & 0.041 & 10.6 \\
\texttt{rand3\_250\_s1} & 61.9 & 32.9 & 5.12 & 0.092 & 2.2 \\
\texttt{rand3\_250\_s3} & 56.2 & 39.2 & 4.51 & 0.083 & 2.1 \\
\texttt{full17} & 97.1 & 2.8 & 0.06 & 0.057 & 8.1 \\
\bottomrule

\end{tabular}
\caption{Outcome of every clause visit at the default deletion schedule.}
\label{tab:visits}
\end{table}

Most visits are skips, 56--97\% (Table~\ref{tab:visits}), yet on
\inst{php11\_10} the skip path takes only about 10\% of the time: the quarter
of visits that search for a new watch take half of the run. A search checks
about 10.6 literals on \inst{php11\_10} but only 2.1 on
\inst{rand3\_250\_s3}, which explains why the two profiles in
Figure~\ref{fig:regions} differ. Between 87\% and 99\% of all visits go to
learnt clauses, so the deletion policy decides how much work propagation has
to do.

One counter disproved an assumption. I expected the first check of each visit,
whether the false literal sits in the first or the second watched position, to
go either way about half the time. It is the first position only about 20\% of
the time (2.5\% on \inst{full17}), because the previous visit usually leaves it
in the second. That branch is therefore easy to predict, and the time around
it is more likely spent waiting for the clause's literals to arrive.

\section{A cycle-level cost model}
\label{sec:costmodel}

Can these counts explain the run time? If every kind of operation has a fixed
cost, then
\begin{equation}
  \text{cycles} \approx a \cdot \text{skips} + b \cdot \text{searches}
    + c \cdot \text{literals checked} + d \cdot \text{conflicts},
\end{equation}
where the searches are the visits that look for a new watch and the conflict
term covers all per-conflict work: analysis, learning, backjumping and the
next decision. I estimated $a$--$d$ by least squares on the relative error over
\numModelRuns{} runs: three families, each under four clause-deletion
schedules, from 0.1 to 26\,s. The schedules change the size of the clause
database and with it the mix of operations, which is what lets the fit tell
the four costs apart.

The result (Table~\ref{tab:coef}) is \numCoefSkip{} cycles per skip,
\numCoefSearch{} per search, \numCoefScan{} per literal checked and
\numCoefConflict{} per conflict, and leaving out any single instance barely
changes it. The model predicts the total cycles of a run with a median error of
\numModelMedianErr\% (Figure~\ref{fig:costmodel}). The costs are also
physically plausible. A skip is cheaper than a single L2 access, which is only
possible because independent visits overlap; a search, about 14 skips, leaves
room for a misprediction plus several dependent loads; and a conflict, about
1.1\,$\mu$s, matches the profile's share for \texttt{analyze()}. The processor
agrees independently: the family dominated by cheap skips runs at 5.0--5.6
instructions per cycle (fullN), those dominated by searches at 1.6--2.3.

On runs it was not fitted on, the model is less reliable. Fitted on pigeonhole
and fullN only, it underpredicts every random 3-SAT run by 15--21\%, partly
because both training families check 8--10 literals per search, so the fit
cannot separate the cost per search from the cost per literal. Fitted on
pigeonhole and random 3-SAT, it overpredicts fullN by 4--47\%, and more so the
more skips a run contains: a skip costs about 9 cycles on those two families
but about 6 on fullN, whose visits are almost all predictable skips over
clauses read largely in order (Section~\ref{sec:memory}). The costs are
accurate within the families they were fitted on, but they depend on the
workload.

\begin{table}[ht]
\centering\small
\begin{tabular}{lrc}
\toprule
cycles per & estimate & leave-one-instance-out range \\
\midrule
% generated by make_data.py - do not edit by hand
visit, clause already satisfied & 6.0 & [5.9, 6.2] \\
visit that searches for a new watch & 86.8 & [85.8, 88.7] \\
literal checked by a search & 2.1 & [1.5, 2.4] \\
conflict & 4232 & [3562, 4642] \\
\bottomrule

\end{tabular}
\caption{Per-operation costs fitted over \numModelRuns{} runs.}
\label{tab:coef}
\end{table}

\begin{figure}[ht]
\centering
\plotfigure{figures/fig-costmodel}
\caption{Measured against predicted cycles for every run (median error
\numModelMedianErr\%). The outlier below the line is \inst{php11\_10} with
clause deletion off, the only run whose working data far exceeds L2.}
\label{fig:costmodel}
\end{figure}

\section{Memory and caches}
\label{sec:memory}

Two results above point to memory: the worst prediction of the cost model was
the run with the largest clause database, and skips are cheapest on fullN. The
second curve in Figure~\ref{fig:memory} (left) explains the latter. Reading
memory in order costs 4--8 cycles at any size, and although \inst{full17} uses
31\,MB, almost twice the L2, it creates its clauses one after another and
mostly visits them in that order. Size alone does not make memory slow; size
combined with a random access order does. Estimated from the data structures
(32 bytes per clause object, 4 bytes per literal in a separate block, 8 bytes
of watch entries), the clause data of \inst{php11\_10} takes about 4\,MB, well
inside L2, so in the normal runs memory is not the bottleneck.

\begin{figure}[ht]
\centering
\plotfigure{figures/fig-memory}
\caption{Left: load latency on this machine, for a random pointer chase and an
in-order walk. Right: extra cycles per clause visit beyond the in-L2 cost
model, as clause padding grows the footprint with the search held fixed. The
dashed line marks the 16\,MB L2.}
\label{fig:memory}
\end{figure}

To isolate memory, I added a setting, \texttt{SAT\_CLAUSE\_PAD}, that reserves
unused space after every clause's literals. The literals are untouched, so the
search is identical (every counter was checked), while the footprint grows up
to \numPadFootprint{} times, from 11 to 226\,MB on \inst{php11\_10}. Comparing
each run with the cost model fitted only on in-L2 runs gives the extra cycles
per visit caused by memory alone. They stay near zero up to an estimated
15\,MB of clause data, rise between 15 and 40\,MB, where the latency curve
leaves L2, and reach about 9 cycles at the largest padding, making each visit
\numPadSlowdown\% more expensive (Figure~\ref{fig:memory}, right). That is far
less than one main-memory access, about 450 cycles, because independent visits
overlap their misses.

\begin{table}[ht]
\centering\small
\begin{tabular}{lrrr}
\toprule
region of \texttt{propagate()} & pad 0 & pad 2048 & change \\
\midrule
% generated by make_data.py - do not edit by hand
replacement search & 17.1 & 19.9 & +2.8 \\
unattributed & 4.0 & 6.1 & +2.1 \\
already satisfied & 3.6 & 5.6 & +2.0 \\
reach the clause & 4.6 & 5.7 & +1.1 \\
watch-list walk & 2.1 & 1.9 & -0.1 \\
move the watch & 0.8 & 0.7 & -0.2 \\
\bottomrule

\end{tabular}
\caption{Cycles per clause visit by code region on \inst{php11\_10}, same
search, unpadded and padded. Only regions that read clause literals grow.}
\label{tab:shift}
\end{table}

A profile of the padded run shows where the cost lands
(Table~\ref{tab:shift}): only in code that reads the padded clause data. The
watch-list walk and the watch move touch only memory the padding did not
change, and they do not slow down, which makes them a control group inside the
same measurement. Padding and a large clause database may leave the cache in
different ways. Padding spreads the same hot data over more memory pages; one
plausible explanation of its cost is that the processor runs out of
address-translation entries (the TLB) rather than out of cache, but these
measurements cannot tell TLB misses from cache misses. With clause deletion
turned off, by contrast, \inst{php11\_10} keeps 166{,}675 learnt clauses and
genuinely reads more data, at about 10 extra cycles per visit and 1.70
instructions per cycle. Without miss counts, four signals agree: the latency
curve, the knee at the L2 size, the falling instructions per cycle, and the
cost appearing only where padded data is read. Together they show that memory
becomes more expensive beyond L2, but not which structure is responsible; a
pointer chase that touches one cache line per page would measure the TLB's
reach and separate the two.

\section{Comparison with MiniSat}
\label{sec:minisat}

A single run made the solver look 12 times faster than
MiniSat~\cite{een2003minisat}, with preprocessing off, on \inst{php11\_10}:
5.4\,s against 65.6\,s. Splitting the run time into conflicts and cost per
conflict shows why. MiniSat needed 42 times more conflicts, because its
learnt-clause limit is sized from the formula: it starts at about 200 clauses
on this 561-clause instance, while this solver keeps up to 23{,}000. Each of
MiniSat's conflicts was 3.5 times cheaper, but that does not show a faster
engine, since with so few clauses there is far less to visit per conflict.
Raising MiniSat's limit to 20{,}000 cut its conflicts tenfold but made each
one 22 times more expensive: at 137\,$\mu$s against this solver's 22\,$\mu$s,
it was then six times \emph{slower} per conflict, while holding 23{,}000 to
74{,}000 clauses. Cost per conflict therefore mostly reflects the size of the
clause database, not the efficiency of the engine. A fair engine comparison
would need the cost per clause visit at an equal database size, which MiniSat
does not report. Clause deletion is a trade-off: fewer clauses make each
conflict cheaper, but there are more conflicts.

The comparison does not survive a luck control. Over 50 shuffled copies of
\inst{php9\_8}, the number of conflicts varies about twofold for the same
formula, and the result depends on the signs of the literals. MiniSat needs
\numGapFlipFull{} times as many conflicts as this solver on copies with sign
flips, but only \numGapNoflipFull{} times as many on copies without
(Figure~\ref{fig:ablation}). Only solvers without restarts are sensitive to
the signs (Figure~\ref{fig:luck}): with its restarts switched off, MiniSat
shows the same twofold shift as this solver. Restarts protect a solver from an
unlucky start~\cite{gomes2000heavytails,luby1993}.

Removing MiniSat's features one at a time closes the gap. Without restarts,
the ratio becomes \numGapFlipNoRestart{} and \numGapNoflipNoRestart{}; without
restarts and clause minimization, it becomes \numGapFlipNoRestartNoMin{} and
\numGapNoflipNoRestartNoMin{}, with confidence intervals that include 1. On
this instance the whole difference in conflicts is these two features, and
even the learnt clauses match: 24.4 and 26.3 literals for this solver, 24.9
and 26.8 for MiniSat without minimization, whose clauses are about 19\%
shorter with it. The
remaining differences, including LBD against activity for ranking learnt
clauses, made no measurable difference here; how many clauses are kept
mattered far more than which ones. Two further observations stay open.
MiniSat's preprocessing increased its conflicts by about 1.3 times, but only
with restarts on, and the original \inst{php11\_10} file is a lucky case for
this solver: 245{,}609 conflicts, against a median of 389{,}126 over shuffled
copies.

\begin{figure}[ht]
\centering
\plotfigure{figures/fig-ablation}
\caption{Conflicts of MiniSat (no preprocessing) divided by conflicts of this
solver, median over 50 shuffled copies of \inst{php9\_8} with 95\% bootstrap
intervals. Removing restarts and clause minimization from MiniSat closes the
gap for both kinds of copies.}
\label{fig:ablation}
\end{figure}

\begin{figure}[ht]
\centering
\plotfigure{figures/fig-luck}
\caption{Conflicts over 50 shuffled copies of \inst{php9\_8}, with and without
sign flips. Only the solvers without restarts are sensitive to the signs.}
\label{fig:luck}
\end{figure}

\section{Which growth is the solver's fault?}
\label{sec:complexity}

Some growth is inherent. On pigeonhole formulas, the median number of
conflicts over shuffled copies grows about 4.7 times per added pigeon, from
724 on \inst{php7\_6} to 389{,}126 on \inst{php11\_10}, while the formula
itself grows only like $n^{3.2}$. No CDCL solver can avoid this: every
resolution proof of the pigeonhole principle is exponentially
long~\cite{haken1985}, and a CDCL run is such a proof. The fullN family, all
$2^n$ clauses over $n$ variables, is the opposite case. Removing any one clause
makes it satisfiable, so every proof must use all $2^n$ clauses, and the
solver needs exactly $2^{n-1}$ conflicts for every $n$ from 13 to 17: work
linear in the size of the input, matching that lower bound up to a constant.

\section{What the measurements predict}
\label{sec:roadmap}

Each planned feature can now be given a measured expectation before it is
built (Table~\ref{tab:roadmap}), so each one becomes a test of the model as
well as an improvement.

\begin{table}[ht]
\centering\small
\begin{tabular}{p{0.24\linewidth}p{0.68\linewidth}}
\toprule
next step & what the measurements predict \\
\midrule
restarts & less sensitivity to variable order and signs (the twofold shift
           should shrink toward MiniSat's 1.1); not necessarily fewer
           conflicts on pigeonhole \\
clause minimization & learnt clauses about 19\% shorter and 13--27\% fewer
           conflicts (Section~\ref{sec:minisat}) \\
blocking literals & saves at most the cost of skips: 8--15\% of cycles on
           pigeonhole, 4--19\% on random 3-SAT, up to 87\% on fullN; more once
           data leaves L2 \\
cheaper search & the largest target, 62--82\% of cycles on pigeonhole and
           random 3-SAT; resuming the scan where the last one
           stopped~\cite{gent2013watched} would shorten it \\
VSIDS heap & no measurable effect at 250 variables; matters at $10^5$ \\
arena storage & removes one pointer indirection per visit, the cost measured
           in Section~\ref{sec:memory} \\
binary DRAT & smaller proofs, less of the
           \numProofOverheadMin--\numProofOverheadMax\% logging overhead \\
\bottomrule
\end{tabular}
\caption{Measured expectations for the planned features.}
\label{tab:roadmap}
\end{table}

\section{Conclusion and open questions}
\label{sec:limitations}

The solver's run time is explained, to within a few percent, by four
per-operation costs, and on a hard combinatorial instance its gap to MiniSat
is exactly two missing features. The most useful lesson was about method: the
first comparison suggested a twelvefold advantage, and only shuffled copies,
paired statistics and removing features one at a time showed what that number
was made of. Three questions remain open: why preprocessing hurts MiniSat only
when restarts are on, which an experiment varying the restart interval could
answer; why proof checking is so slow on fullN; and whether ranking clauses by
LBD pays off on structured industrial instances, which the benchmark set here
does not contain.


\FloatBarrier   % every float is placed before the references begin

\bibliographystyle{plain}
\bibliography{refs}

\appendix
\section{Reproducing the results}
\label{app:repro}

All commands run from the repository root after \texttt{cd src \&\& make
release \&\& make bench-gen}. Each script writes a CSV whose header records its
configuration, and \texttt{docs/report/make\_data.py} turns those files into
every figure, table and number in this report.

\begin{center}\small
\begin{tabular}{ll}
\toprule
result & command \\
\midrule
regression and proofs & \texttt{cd src \&\& make test} \\
fuzzing & \texttt{cd src \&\& make fuzz \&\& ./fuzz 15000} \\
profile (Sec.~\ref{sec:profile}) & \texttt{sample <pid> 4 1 -file php.sample.txt} \\
cost model (Sec.~\ref{sec:costmodel}) & \texttt{python3 tests/cost\_model.py collect} \\
memory sweep (Sec.~\ref{sec:memory}) & \texttt{python3 tests/memory\_sweep.py run} \\
latency ladder & \texttt{clang -O2 tests/latency\_ladder.c \&\& ./a.out} \\
proof cost (Sec.~\ref{sec:solver}) & \texttt{python3 tests/proof\_cost.py --lrat} \\
proof complexity (Sec.~\ref{sec:complexity}) & \texttt{python3 tests/proof\_complexity.py} \\
report data & \texttt{python3 docs/report/make\_data.py} \\
\bottomrule
\end{tabular}
\end{center}

The MiniSat comparison (Section~\ref{sec:minisat}) runs
\texttt{tests/luck\_test.py tests/bench/php9\_8.cnf --seeds 1-50} six times:
with and without \texttt{--no-flip}, each with default MiniSat, with
\texttt{--minisat-flags "-rfirst=2147483647" --tag norestart}, and with
\texttt{--minisat-flags "-rfirst=2147483647 -ccmin-mode=0" --tag
norestart-nomin}.


\end{document}
